\documentclass[11pt]{article}
\usepackage[margin=1in]{geometry}
\usepackage{amsmath,amssymb}
\usepackage{graphicx}
\usepackage{hyperref}

\title{Active Inference-Inspired Delegated Procurement: Quantifying Epistemic Uncertainty in High-Stakes Markets}
\author{Darsh Garg, Erina Karati, \& Muhil Ramesh\\University of Minnesota -- CSCI 5512 (AI II)\\Professor Paul Schrater}
\date{March 2026}

\begin{document}

\maketitle

\begin{abstract}
We study \textbf{autonomous high-stakes procurement} in stochastic markets, where an agent must decide whether to purchase an item, ask a clarifying question, continue searching, or wait for a better opportunity. Current LLM-based shopping agents can rank items and generate fluent rationales, but they do not naturally expose calibrated, decision-relevant uncertainty for high-value transactions. We propose a \textbf{model-based delegation engine}, inspired by active inference, that maintains explicit Bayesian beliefs over latent user preferences and combines \textbf{epistemic uncertainty}, \textbf{stochastic market dynamics}, and \textbf{minimax regret} to determine when autonomous purchase is justified. The agent purchases only when worst-case regret is small and epistemic uncertainty is sufficiently low; otherwise it queries the user or gathers more information. To keep the project feasible within the course timeline, we evaluate the method in a stochastic simulator built from a product catalog with simulated price and availability changes.
\end{abstract}

\section{Introduction and Motivation}

In an emerging agent-mediated economy, autonomous buyers may eventually negotiate and execute high-value transactions, for example purchasing a \$1{,}000 laptop or enterprise subscription, on behalf of human users. While large language models (LLMs) and web-scale recommenders can rank items and generate fluent rationales, they are fundamentally black-box function approximators and do not naturally expose calibrated uncertainty over future user satisfaction or market outcomes.\cite{charpentier2022disentangling} As a result, they are not by themselves sufficient for deciding when it is appropriate to commit to a purchase, particularly in volatile, partially observable markets.

We address the problem of \emph{autonomous high-stakes procurement} under uncertainty. Our central design principle is that a buyer agent must (i) be model-based, with explicit probabilistic representations of user preferences and environment dynamics; (ii) quantify its own epistemic uncertainty over the consequences of a purchase; and (iii) optimize a \emph{minimax regret} objective rather than a single point estimate of expected value.\cite{rigter2020minimax,regretDecision1982} When the worst-case regret is small and epistemic uncertainty is low, the agent may autonomously execute a transaction. When regret or epistemic uncertainty is large, the agent should instead defer to the user, either by asking a clarifying question or by abstaining from purchase if no safe decision emerges by the episode horizon. This formulation directly ties to course themes in Bayesian inference, decision making under uncertainty, and stochastic control.\cite{charpentier2022disentangling}

Our proposal is also motivated by the observation that autonomous procurement is not simply a ranking problem. The core challenge is deciding \emph{when an agent should act at all} under incomplete knowledge of both user preferences and market dynamics. In this sense, the project is fundamentally about uncertainty-aware delegation: the agent must trade off immediate reward, information gathering, and the cost of a potentially regrettable purchase.

\section{Problem Statement}

We formalize autonomous procurement as a stochastic control problem with partial observability. Time is discrete, $t = 0,1,\dots,T$. At each step $t$, the agent observes an information state $o_t$ derived from an underlying environment state $s_t \in \mathcal{S}$, selects an action $a_t$ from a finite action set $\mathcal{A}$, and transitions to $s_{t+1}$ according to unknown stochastic dynamics. The episode terminates when the agent executes a \textsc{Purchase} action or the horizon $T$ is reached. In the latter case, abstention is modeled implicitly as ending the episode without purchase when no action satisfies the safety criterion.

\subsection{Action Space and Environment}

The action set is
\begin{equation}
\mathcal{A} = \{\textsc{Search},\, \textsc{Purchase},\, \textsc{QueryUser},\, \textsc{Wait}\}.
\end{equation}

The environment includes a catalog of items $i \in \mathcal{I}$, each with a feature vector $x_i \in \mathbb{R}^d$ capturing price, rating, category, and other attributes. Item availability and effective price evolve stochastically due to stock-out events, discounts, and competing buyers.

We assume a parametric environment model
\begin{equation}
P(s_{t+1} \mid s_t, a_t, \theta_E),
\end{equation}
where $\theta_E$ parameterizes price and availability dynamics. The true environment parameters are unknown; the agent operates with an approximate model that captures key stochastic structure, such as the probability of stock-out as a function of demand and rating.

The \textsc{Wait} action allows the agent to defer execution when the expected temporal value of waiting is positive, meaning the agent anticipates that the price or availability model may transition into a more favorable state, such as after a predicted discount or restock, before committing to purchase. Concretely, the agent must decide when to purchase, when to ask a clarifying question, when to continue gathering information, and when to defer action because uncertainty remains too high.

\subsection{User Utility and Minimax Regret}

User preferences over item features are represented by a latent utility function
\begin{equation}
U_{\text{user}}(x; \theta_U),
\end{equation}
where $\theta_U$ is a vector of unknown preference parameters, for example trade-offs between price, quality, and brand. The agent does not observe $\theta_U$ directly and must infer it from sparse signals, such as explicit ratings, binary feedback, or answers to targeted queries.

For a given environment state and unknown $\theta_U$, an optimal clairvoyant decision-maker over the currently available set could choose the item
\begin{equation}
i_t^\star(\theta_U) = \arg\max_{i \in \mathcal{I}_t} U_{\text{user}}(x_i; \theta_U).
\end{equation}
If the agent instead selects item $i$, the \emph{regret} under preference parameters $\theta_U$ is
\begin{equation}
R(i; \theta_U) = U_{\text{user}}(x_{i_t^\star(\theta_U)}; \theta_U) - U_{\text{user}}(x_i; \theta_U).
\end{equation}

Because $\theta_U$ is unknown, the agent maintains a belief distribution $b_t(\theta_U)$ and evaluates \emph{minimax} or worst-case regret over a credible set of parameters $\Theta_t$ induced by this belief, for example a high-probability ellipsoid around the posterior mean.\cite{rigter2020minimax} The decision task is to choose a policy that, when it does purchase, keeps worst-case regret below an acceptable threshold.\cite{regretDecision1982}

\section{Background and Related Work}

\subsection{Web Agents and Delegated Shopping}

Recent work on language-grounded web agents has introduced environments where agents navigate product listings and web pages to fulfill natural language shopping instructions.\cite{yao2022webshop} These systems often leverage LLMs to generate actions and rationales over semi-structured web content.\cite{yao2022webshop} While they demonstrate strong performance on benchmark tasks, they typically optimize success rates or rewards without explicit uncertainty quantification or regret guarantees. As a result, they do not provide principled mechanisms for deciding when to defer to a human in high-stakes scenarios.

\subsection{Uncertainty-Aware Recommendation and Epistemic Uncertainty}

Uncertainty-aware recommendation systems model user-item interactions probabilistically, estimating predictive distributions rather than point predictions. Bayesian and ensemble-based approaches distinguish between aleatoric (inherent noise) and epistemic (model) uncertainty, using the latter to improve robustness and to avoid overconfident recommendations on sparse or out-of-distribution data.\cite{bdecf2025} In this work, we adopt a Bayesian formulation of user preferences and treat epistemic variance as a first-class quantity that drives both minimax regret and a purchase safety gate.\cite{charpentier2022disentangling}

\subsection{Minimax Regret and Robust Decision Making}

Minimax regret is a classical criterion in robust decision theory: instead of maximizing expected utility under a single believed model, a decision-maker chooses an action that minimizes the maximum regret over a set of plausible models.\cite{regretDecision1982} This yields policies that are conservative with respect to model misspecification and uncertainty about underlying parameters. In stochastic environments, minimax regret can be combined with Bayesian belief updates to balance robustness and data efficiency.\cite{rigter2020minimax} Our agent uses minimax regret over a belief-defined confidence set to decide when autonomous purchasing is acceptable.

\subsection{Active Inference and Information-Seeking Action}

Active inference provides a complementary perspective in which an agent maintains a generative model over hidden states, updates beliefs as evidence arrives, and selects actions that jointly pursue preferred outcomes and reduce uncertainty.\cite{bogacz2017tutorial,smith2022active} This framework is particularly relevant to delegated procurement because actions such as \textsc{QueryUser} and \textsc{Search} are not merely fallback behaviors; they are \emph{epistemic actions} that improve the agent's knowledge of user preferences or market conditions before committing to a purchase. This makes active inference a useful conceptual lens for interpreting uncertainty-aware delegation.

\subsection{Positioning}

Our contribution lies at the intersection of delegated web agents, uncertainty-aware recommendation, robust decision making, and active inference.\cite{yao2022webshop,bdecf2025,rigter2020minimax,bogacz2017tutorial} Unlike standard LLM-based shopping agents, our model-based delegation engine maintains explicit probabilistic beliefs over user preferences and environment dynamics. Unlike typical recommenders that output ranked lists, our agent must decide \emph{whether} to act autonomously at all, using minimax regret and epistemic variance as safety metrics. This combination yields an uncertainty-aware decision layer tailored to delegated purchasing under stochasticity.

\section{Proposed Methodology}

We now describe our modeling choices and decision rule in more detail. Throughout, stochasticity and probabilistic reasoning are central.

While classical active inference is often formulated through variational free energy minimization, our implementation is better understood as \emph{active inference-inspired}: the agent uses Bayesian belief updates and information-seeking actions, especially \textsc{QueryUser}, to reduce epistemic uncertainty before purchase.

\subsection{User Preference Model and Belief State}

We represent each item by a feature vector $x \in \mathbb{R}^d$, for example normalized price, quality score, and brand indicator. The user's latent utility is modeled as
\begin{equation}
U_{\text{user}}(x; \theta_U) = \theta_U^\top x + \epsilon, \quad \epsilon \sim \mathcal{N}(0, \sigma_\epsilon^2),
\end{equation}
with a Gaussian prior over preferences,
\begin{equation}
\theta_U \sim \mathcal{N}(m_0, S_0).
\end{equation}

As the agent observes user signals, such as ``I care more about durability than style'' or ratings on probed items, it updates a posterior $p(\theta_U \mid \mathcal{D}_t)$, which remains Gaussian under linear-Gaussian assumptions. The belief state is summarized by $(m_t, S_t)$.

For a candidate item $x$, the predictive distribution over utility is
\begin{equation}
U_{\text{user}}(x) \mid \mathcal{D}_t \sim \mathcal{N}\big(m_t^\top x,\; x^\top S_t x + \sigma_\epsilon^2\big),
\end{equation}
so the agent has both an expected utility and an uncertainty estimate for each item.

We distinguish \emph{epistemic uncertainty}, given by $x^\top S_t x$, from total predictive uncertainty, given by $x^\top S_t x + \sigma_\epsilon^2$. Since epistemic uncertainty is reducible through additional evidence, it is the primary uncertainty term used in the purchase safety gate.

\subsection{Environment Stochasticity and Digital Twin}

To model market volatility, we define a simplified stochastic process over availability and effective price. Let $z_{i,t} \in \{0,1\}$ indicate whether item $i$ is in stock at time $t$. We specify a stock-out probability
\begin{equation}
\mathbb{P}(z_{i,t+1} = 0 \mid z_{i,t} = 1) = \alpha_i,
\end{equation}
where $\alpha_i$ is a function of item attributes such as demand, rating, and category. Similarly, we may model stochastic discounts as a random process on price.

In our implementation, this environment model serves as a stochastic simulator: before committing to a final purchase decision, the agent runs Monte Carlo rollouts to estimate distributions over future outcomes under candidate policies. This explicitly incorporates environmental stochasticity into the decision rule.

\subsection{Minimax Regret Objective}

Given a belief $p(\theta_U \mid \mathcal{D}_t)$, we define a high-probability confidence set
\begin{equation}
\Theta_t = \{\theta_U : (\theta_U - m_t)^\top S_t^{-1} (\theta_U - m_t) \leq c_t\},
\end{equation}
for some radius $c_t$ determined by a chosen confidence level. For any item $i$, the worst-case regret over $\Theta_t$ is
\begin{equation}
\bar{R}_t(i) = \sup_{\theta_U \in \Theta_t} \Big[ \max_{j \in \mathcal{I}_t} U_{\text{user}}(x_j; \theta_U) - U_{\text{user}}(x_i; \theta_U) \Big],
\end{equation}
where $\mathcal{I}_t$ is the set of currently available items. The agent seeks items with small $\bar{R}_t(i)$; in practice we approximate $\bar{R}_t(i)$ using linear bounds or sampling over $\theta_U$ from the confidence set.\cite{rigter2020minimax}

\subsection{Safety Gate and Delegation Decision}

At a decision point, for each candidate item $i$ the agent computes
\begin{align}
\mu_t(i) &= m_t^\top x_i,\\
\sigma_{\text{epi},t}^2(i) &= x_i^\top S_t x_i,\\
\sigma_{\text{pred},t}^2(i) &= x_i^\top S_t x_i + \sigma_\epsilon^2,\\
\bar{R}_t(i) &\approx \text{approximate worst-case regret over } \Theta_t.
\end{align}

We introduce two thresholds:
\begin{itemize}
    \item $\epsilon_{\text{reg}}$: maximum tolerable worst-case regret.
    \item $\epsilon_{\text{epi}}$: maximum tolerable epistemic uncertainty.
\end{itemize}

The \emph{autonomous purchase condition} is
\begin{equation}
\bar{R}_t(i) \leq \epsilon_{\text{reg}} \quad \text{and} \quad \sigma_{\text{epi},t}^2(i) \leq \epsilon_{\text{epi}}.
\end{equation}

If there exists an item satisfying this condition and with sufficiently high mean utility, the agent may execute \textsc{Purchase}. Otherwise, the agent gathers more information by choosing \textsc{QueryUser}, \textsc{Search}, or \textsc{Wait}. We still track predictive uncertainty $\sigma_{\text{pred},t}^2(i)$ for analysis, but the safety gate is driven primarily by epistemic uncertainty because it is reducible through additional evidence.

\subsection{Connection to Active Inference}

Our formulation is also compatible with an active inference perspective, in which the agent maintains a generative model over hidden user preferences and market dynamics, updates beliefs as new evidence arrives, and selects actions that reduce uncertainty while moving toward preferred outcomes.\cite{bogacz2017tutorial,smith2022active} In the MVP, the maintained belief state is primarily over latent user preference parameters $\theta_U$; market dynamics are stochastic but treated as a fixed simulator model rather than jointly inferred hidden variables. In this framing, \textsc{QueryUser} and \textsc{Search} are not merely fallback actions; they are epistemic actions that reduce uncertainty in the agent's belief state. Operationally, \textsc{QueryUser} acts as a targeted preference-elicitation step over the latent utility model. This aligns with active inference accounts in which action selection trades off goal-directed behavior with information gain under a probabilistic generative model.\cite{smith2022active}

In our setting, this connection appears operationally through the expected information gain term
\begin{equation}
IG(a) = \mathcal{H}(b_t) - \mathbb{E}[\mathcal{H}(b_{t+1})],
\end{equation}
which measures how much an action is expected to reduce uncertainty over the user-preference belief state. Thus, although our implementation is presented in minimax-regret form, it can also be interpreted as an active inference-inspired delegation agent that prefers autonomous purchase only when both expected utility is high and epistemic uncertainty is sufficiently low.

\subsection{Pseudocode for the Delegation Engine}

Below we sketch the core decision logic.

\medskip
\noindent\textbf{Algorithm 1: Model-Based Delegation Engine}
\begin{enumerate}
    \item Input: belief $(m_t, S_t)$, available items $\mathcal{I}_t$, thresholds $\epsilon_{\text{reg}}, \epsilon_{\text{epi}}$, utility threshold $\tau_{\text{util}}$.
    \item For each $i \in \mathcal{I}_t$:
    \begin{enumerate}
        \item Compute $\mu_t(i)$, $\sigma_{\text{epi},t}^2(i)$, and approximate worst-case regret $\bar{R}_t(i)$ over $\Theta_t$.
    \end{enumerate}
    \item Define the safe candidate set
    \[
    \mathcal{C}_t = \left\{ i \in \mathcal{I}_t \,:\, \mu_t(i) \geq \tau_{\text{util}},\ \bar{R}_t(i) \leq \epsilon_{\text{reg}},\ \sigma_{\text{epi},t}^2(i) \leq \epsilon_{\text{epi}} \right\}.
    \]
    \item If $\mathcal{C}_t \neq \emptyset$:
    \begin{itemize}
        \item Return \textsc{Purchase}$(i^\star)$ where $i^\star = \arg\max_{i \in \mathcal{C}_t} \mu_t(i)$.
    \end{itemize}
    \item Else:
    \begin{itemize}
        \item Estimate $IG(\textsc{QueryUser})$ as the expected reduction in entropy of the posterior over $\theta_U$ after a user answer.
        \item Estimate the one-step value of \textsc{Wait} using Monte Carlo rollouts under stochastic price and availability dynamics.
        \item If $IG(\textsc{QueryUser})$ is largest, return \textsc{QueryUser}.
        \item Else if the expected value of \textsc{Wait} exceeds the expected value of continued search, return \textsc{Wait}.
        \item Else return \textsc{Search}.
    \end{itemize}
\end{enumerate}

This algorithm is inherently stochastic: it relies on probabilistic beliefs over $\theta_U$, predictive distributions over utilities, randomized market transitions, Monte Carlo rollouts in the simulator, and a confidence set $\Theta_t$ that evolves under Bayesian updates.

\section{Data and Environment}

To keep the project feasible while still demonstrating the full stochastic decision pipeline, we adopt a two-level environment plan.

\paragraph{Static CSV Simulator (Core MVP).} The primary environment is a static catalog stored as a CSV with approximately 1{,}000 products. Each row contains curated attributes: price, rating, category, and derived quality metrics. We simulate availability and price transitions using the stochastic stock-out model and, optionally, random discount processes. The agent interacts with this environment entirely in simulation, enabling rapid experimentation. The catalog will be derived from a curated public dataset or a controlled synthetic dataset, depending on data availability and preprocessing time.

\paragraph{Stochastic Simulator Interface.} In the static simulator, the planning model and the environment coincide. If time permits, we will connect the same delegation engine to a more realistic web environment, where the simulator is used for Monte Carlo planning prior to issuing real actions.

Synthetic user personas, for example a ``budget shopper'' or ``quality maximizer,'' are defined by different priors $(m_0, S_0)$ over $\theta_U$, providing hidden preference structure while avoiding privacy issues.

\paragraph{Feasibility and Scope Control.} To ensure the project remains tractable within the course timeline, we intentionally begin with a controlled simulator rather than full web deployment. This allows us to isolate the core scientific question, namely whether explicit uncertainty quantification and minimax-regret planning improve delegated purchasing decisions under stochasticity, before introducing additional engineering complexity. If integration with a live web environment proves too time-consuming, the static simulator will still fully support the proposed stochastic modeling, Bayesian updates, Monte Carlo evaluation, and baseline comparisons.

\section{Implementation Plan and Timeline}

We assume a three-person team, with work being organized around modular components.

\subsection{Roles}

\begin{itemize}
    \item \textbf{Member 1}: Implement the Bayesian user preference model, posterior updates, and confidence set construction.
    \item \textbf{Member 2}: Implement the stochastic stock-out and price dynamics, simulator logic, and feature processing from the CSV catalog.
    \item \textbf{Member 3}: Implement the minimax-regret approximation, delegation engine, experiment pipeline, and evaluation scripts.
\end{itemize}

\subsection{Timeline}

\paragraph{Weeks 1--2: Foundations.} Finalize the problem specification and metrics. Implement a prototype Bayesian user model and generate synthetic personas. Load a toy CSV dataset and basic feature extraction.

\paragraph{Week 3: Environment Stochasticity.} Implement the stock-out and price dynamics, environment step function, and item sampling procedures. Verify that the simulator produces realistic variability across episodes.

\paragraph{Week 4: Delegation Engine.} Implement the minimax-regret calculation using sampling or linear bounds and the epistemic-uncertainty safety gate. Integrate thresholds and decision logic into a clean API.

\paragraph{Week 5: Integration and Debugging.} Connect the user model, environment, and planner in a single loop. Run smoke tests demonstrating Buy versus Query decisions under different priors and thresholds.

\paragraph{Week 6: Evaluation and Analysis.} Run experiments across personas and stochastic seeds. Collect metrics, analyze trade-offs between regret, number of queries, no-purchase rate, and purchase rate, and produce plots and tables for the final report.

\section{Evaluation and Success Criteria}

All evaluation will be conducted across multiple stochastic seeds, sampled user preference realizations, and randomized market transitions so that performance reflects robustness under uncertainty rather than success on a single deterministic scenario.

\begin{itemize}
    \item \textbf{Average Regret}: Expected regret of the purchased item relative to the clairvoyant optimum, averaged over episodes.
    \item \textbf{Worst-Case Regret}: Empirical approximation of worst-case regret over sampled $\theta_U$ from the confidence set, highlighting robustness.
    \item \textbf{Delegation / No-Purchase Rate}: Frequency of \textsc{QueryUser} actions or episodes that end without purchase; we expect this to increase when priors are broad and thresholds are strict.
    \item \textbf{Regret Exceedance Rate}: Frequency with which the realized regret of a purchased item exceeds the target tolerance under sampled true preferences, measuring whether the safety gate is well-calibrated under stochasticity.
    \item \textbf{Persona Adaptation}: Differences in behavior across synthetic personas, demonstrating that changing priors alone changes the agent's strategy without retraining.
\end{itemize}

We will compare the model-based minimax-regret agent against a baseline that selects items solely by posterior mean utility $\mu_t(i)$, without explicit regret or uncertainty constraints. We hypothesize that the baseline will either incur higher regret or require ad hoc heuristics to avoid poor purchases, whereas our approach uses principled probabilistic criteria.

We will also perform ablations that remove (i) posterior uncertainty estimates, (ii) the minimax-regret term, and (iii) stochastic market dynamics, in order to isolate which components are responsible for improvements in safety and robustness.

We will consider the project successful if, relative to the posterior-mean baseline, the proposed agent reduces empirical regret and regret exceedance rate while maintaining a reasonable purchase rate across stochastic seeds and user personas.

\section{Conclusion}

We propose a model-based delegation engine for autonomous high-stakes procurement in stochastic markets. By explicitly modeling latent user preferences and stochastic market dynamics, our system treats uncertainty, regret, and delegation as first-class decision variables rather than after-the-fact heuristics. By quantifying epistemic uncertainty and optimizing a minimax-regret objective, the agent provides a principled uncertainty-aware decision layer for deciding when to buy, when to ask, and when to defer.\cite{charpentier2022disentangling,rigter2020minimax} This directly addresses the course requirement that the project incorporate stochasticity into the core decision-making process.

\begin{thebibliography}{9}

\bibitem{yao2022webshop}
Shunyu Yao, Howard Chen, John Yang, and Karthik Narasimhan.
\newblock WebShop: Towards Scalable Real-World Web Interaction with Grounded Language Agents.
\newblock \emph{Advances in Neural Information Processing Systems (NeurIPS)}, 2022.

\bibitem{bdecf2025}
Radin Cheraghi, Amir Mohammad Mahfoozi, Sepehr Zolfaghari, Mohammadshayan Shabani, Maryam Ramezani, and Hamid R. Rabiee.
\newblock Epistemic Uncertainty-aware Recommendation Systems via Bayesian Deep Ensemble Learning.
\newblock \emph{arXiv preprint arXiv:2504.10753}, 2025.

\bibitem{charpentier2022disentangling}
Bertrand Charpentier, Ransalu Senanayake, Mykel J. Kochenderfer, and Stephan G{\"u}nnemann.
\newblock Disentangling Epistemic and Aleatoric Uncertainty in Reinforcement Learning.
\newblock \emph{arXiv preprint arXiv:2206.01558}, 2022.

\bibitem{rigter2020minimax}
Marc Rigter, Bruno Lacerda, and Nick Hawes.
\newblock Minimax Regret Optimisation for Robust Planning in Uncertain Markov Decision Processes.
\newblock In \emph{Proceedings of the Conference on Robot Learning (CoRL)}, 2020.

\bibitem{regretDecision1982}
David E. Bell.
\newblock Regret in Decision Making under Uncertainty.
\newblock \emph{Operations Research}, 30(5):961--981, 1982.

\bibitem{bogacz2017tutorial}
Rafal Bogacz.
\newblock A tutorial on the free-energy framework for modelling perception and learning.
\newblock \emph{Journal of Mathematical Psychology}, 76:198--211, 2017.

\bibitem{smith2022active}
Ryan Smith, Thomas Parr, and Karl J. Friston.
\newblock A step-by-step tutorial on active inference and its application to empirical data.
\newblock \emph{Journal of Mathematical Psychology}, 107:102632, 2022.

\end{thebibliography}

\end{document}
